\documentclass[12pt]{article}
\usepackage[utf8]{inputenc}
\usepackage[T1]{fontenc}
\usepackage{hyperref}
\title{Lusitanian Language Resources for K-12 Curriculum Development}
\author{Compiled by Claude Sonnet 5.5 (claude-sonnet-5-5) for Andr\'as Kornai}
\date{}
\begin{document}
\maketitle

\section{Introduction}
Lusitanian (ISO 639-3 xls; Glottocode lusi1235) is an extinct Indo-European language of the western Iberian Peninsula. It has no native speakers. The search results describe a very small corpus of about six inscriptions, all religious in content and all written in the Latin alphabet: Arroyo de la Luz I--III, Lamas de Moledo, Cabeço das Fráguas and Ribeira da Venda (found in 2008). The Arroyo de la Luz texts survive only through copies. Dating of the inscriptions is uncertain. Because no teaching materials exist, a school programme can only be a heritage-awareness and inquiry unit built on the scholarly literature. It cannot teach the language for communication. All sources below are in English, Spanish or Portuguese, and none are in Lusitanian.

\section{Primary Educational Resources}
I found no primers, textbooks or school programmes for Lusitanian. Teachers would need to adapt the scholarly or popular materials listed in the other sections.

\section{Linguistic and Reference Materials}
\begin{itemize}
\item General overview of the Lusitanian language and its inscriptions: \url{https://wikipedia.com/wiki/ISO_639:xls}
\item Book chapter on religious inscriptions in north-western Roman Hispania (Presses universitaires de Nouvelle-Aquitaine): \url{https://una-editions.fr/writing-to-the-gods-in-north-western-roman-hispania/}
\item Study by Stifter and colleagues on the Lusitanian material (Maynooth University repository): \url{https://mural.maynoothuniversity.ie/11398/1/Stifter}
\item Spanish-language study of Lusitanian epigraphy (Universidad Complutense, Docta): \url{https://docta.ucm.es/rest/api/core/bitstreams/5bcef4ad-a05e-4e0d-9d8d-1205cbabca97/content}
\item Spanish-language thesis text (Universidad de Zaragoza) that treats the six inscriptions as the whole Lusitanian corpus: \url{https://zaguan.unizar.es/record/70251/files/texto_completo.pdf}
\item Masaryk University digital library PDF: \url{https://digilib.phil.muni.cz/_flysystem/fedora/pdf/114048.pdf}
\item Academia.edu papers returned by the search: \url{https://www.academia.edu/9563554}, \url{https://www.academia.edu/31989410}, \url{https://www.academia.edu/193620}
\end{itemize}
The search summaries only gave the gist of these items, and some were OCR-garbled. I have not verified their authorship or content in detail. Check the primary editions (for example the Untermann and Villar corpora) before relying on specifics.

\section{Religious and Community Resources}
The inscriptions are themselves religious texts: dedications and offerings naming local deities and animals, notably at Cabeço das Fráguas. I found no living community organisations and no scripture translations.

\section{Digital and Audio Resources}
I found no audio, video or software resources for Lusitanian. The inscriptions are in the Latin alphabet, so images and transcriptions can be shown in class.

\section{Curriculum Development Context}
There is no speaker community and no pedagogical tradition. A realistic K-12 use is a short cross-curricular unit on the pre-Roman languages of Iberia. Students could read the Latin-alphabet inscriptions, discuss why the corpus is so small, and look at the Cabeço das Fráguas text, the most analysed, with its names of animals and deities. Any ``revived'' vocabulary would be reconstruction and not attested usage. Teachers should say so.
\end{document}